\documentclass[a4paper]{jsarticle}

\usepackage[dvipdfmx]{graphicx}
\usepackage[dvipdfmx]{hyperref}
\usepackage{amsmath}
\usepackage{url}

\graphicspath{{figures/}}

\title{過去就活データを活用した就活支援チャットボットの開発・運用と評価}
\author{鈴木源吾ゼミ \\ 20123005 糸賀 透也}
\date{}

\begin{document}

\maketitle

\tableofcontents
\newpage

% ============================================================
\section{はじめに}
\label{sec:intro}
% ============================================================

\subsection{背景と課題}
\label{sec:background}
% TODO: キャリアセンターに過去の就活情報が蓄積されているにもかかわらず、
%       学生が利用できる窓口はキャリアセンターへの問い合わせのみである点を述べる
% TODO: キャリアセンター職員が情報の所在を把握しきれず、情報取得に時間がかかる点を述べる

\subsection{研究目的}
\label{sec:purpose}
% TODO: 目的①：過去の就活情報を効率的に活用できるアプリを開発することを述べる
% TODO: 目的②：試験を実施してアプリの効果測定を行うことを述べる

% ============================================================
\section{キャリアナビAI の開発}
\label{sec:development}
% ============================================================

\subsection{システム概要}
\label{sec:overview}
% TODO: RAGアーキテクチャの流れ（質問送信→ベクトル検索→データ取得→AI回答生成→表示）を説明する
% FIG: システムの処理フロー図（スライド p.5 の概念図に相当）

\subsection{技術スタック}
\label{sec:techstack}
% TODO: 各レイヤー（フロントエンド・バックエンド・実行環境・認証・DB・ストレージ・生成AI・埋め込みモデル・データ処理・開発デプロイ）の
%       採用技術と役割を表形式でまとめる
% FIG: 技術スタック一覧表（スライド p.6 の表に相当）

\subsection{システム構成}
\label{sec:architecture}
% TODO: AWS 上の構成（Lambda・S3・DynamoDB・Bedrock）と
%       GitHub / GitHub Actions によるデプロイ経路を説明する
% TODO: 教職員ユーザーのみファイル管理可能というロール設計にも言及する
% FIG: システム構成図（スライド p.7 の構成図に相当）

% ============================================================
\section{試験と評価}
\label{sec:evaluation}
% ============================================================

\subsection{試験概要}
\label{sec:trial-overview}
% TODO: 倫理委員会申請を経た上で試験を実施したことに言及する
% TODO: 対象（情報学部3年生10名・教職員6名）、期間（8/31〜9/14、約3週間）、目的を述べる
% TODO: 試験後アンケートの実施についても説明する

\subsection{試験結果}
\label{sec:trial-result}
% TODO: アンケート結果を記載する（定量・定性データ）
% TODO: 要確認：アンケートの設問項目・回答結果の数値
% FIG: アンケート結果の集計グラフ・表（データ取得後に追加）

\subsection{考察}
\label{sec:discussion}
% TODO: 試験結果から読み取れるアプリの有効性・課題を考察する
% TODO: 学生への影響・教職員の業務効率への影響それぞれについて述べる

% ============================================================
\section{機能改善と今後の課題}
\label{sec:future}
% ============================================================
% TODO: アンケート結果を基にした機能追加・改善の内容を述べる
% TODO: 本格導入に向けた検討（ECS移行・保守管理体制の整備等）を述べる
% TODO: 要確認：具体的な改善内容はアンケート分析後に確定

% ============================================================
\section{おわりに}
\label{sec:conclusion}
% ============================================================
% TODO: 本研究の成果をまとめる
% TODO: 本格導入・将来展望について述べる

\bibliographystyle{jplain}
\bibliography{references}

\end{document}
